\section{Repository audit and the precise missing bridge}

\subsection{Pinned source tree}

This report was prepared against the public ProveIt repository at commit
\RepoCommit.  The relevant subtree distinguishes two kinds of material.
The \code{ConciseIPRA} directory is a Lean development whose kernel-checked
results include the ordinary first-order presentation of iPRA, a compact native
scalar calculus, and translations between them.  The neighboring
\code{docs/proofs} directory is a research archive.  At the pinned snapshot its
index records seventy-two delivered packets in algebra, analysis, analytic and
Diophantine number theory, combinatorics, foundations, geometry, probability,
and topology.  The index explicitly warns that these packets were filed as
delivered: their Lean files were not systematically built, and their claims
must not be confused with the checked core \cite{ProveItRoot,PacketsIndex}.

This distinction determines the standard of evidence used here.  A theorem in
this report is called \emph{embedded} only when there is a specified route to a
value of the repository's inductive derivability relation
\[
        \Prf\,\PRAAx\,\Gamma\,\varphi.
\]
A host-language proof about natural numbers, a successful exact-arithmetic
program, or a semantic theorem about a Lean structure is useful evidence, but
it is not itself such a value.

\subsection{What the existing frameworks already accomplish}

The foundation packets are unusually convergent.  Although they use different
terminology, they repeatedly discover the same useful replacements for
classical infinitary objects.

\begin{longtable}{p{0.21\textwidth}p{0.34\textwidth}p{0.37\textwidth}}
\toprule
Classical idiom & Finitary replacement & Recurring repository contribution \\
\midrule
\endfirsthead
\toprule
Classical idiom & Finitary replacement & Recurring repository contribution \\
\midrule
\endhead
A completed infinite object & finite observations indexed by precision &
certified finitary presentations, observation towers, finite names \\
A canonical approximate witness & a finite cloud of verified answers &
relational certificates, root clusters, stable observables \\
Continuity ``by compactness'' & an explicit modulus plus a finite net &
finite observation and uniform-certificate frameworks \\
A long algorithmic argument & a bounded trace with a quantifier-free invariant &
finite-demand, local--global, and compiler packets \\
Nested numerical constructions & backward propagation of required precision &
finite-demand and contract-DAG architectures \\
A partial operation & a total program plus a positive rational guard &
guarded inverses, pivots, charts, and branch rules \\
Finiteness of a Diophantine set & a height producer plus bounded enumeration &
height/list interfaces and finite witness packages \\
Equality of abstract objects & observational equivalence at every finite level &
setoid and relational back-and-forth methods \\
\bottomrule
\end{longtable}

The most developed unified packet proves human-level theorems for bounded
Booleanization, finite trace reflection, relational viability, shrinking
candidate clouds, guarded arithmetic, finite nets, and height/list conversion.
Its theorem ledger nevertheless marks the generic \code{Prf PRAAx} compiler as
an open engineering task.  The subsequent contract-DAG packet strengthens the
architecture with reverse demand propagation and a global trace, but again
leaves the generic object-language theorem unimplemented.  This is the point
at which the present work begins: it treats those semantic ideas as prior art
and tightens the exact interface with the checked \code{ConciseIPRA} API.

\subsection{The checked backend that can be reused}

The file \code{PR.lean} supplies precisely the syntax needed for the target.
It defines primitive-recursive symbols
\[
 \PR(k),
\]
terms \(\PTm(n)\), the axiom family \(\PRAAx\), and the derived rule
\code{PRA.qfInd}.  Primitive recursion is represented by
\code{PR.prec}; its zero and successor equations are available as theorems
\code{PRA.precZero} and \code{PRA.precSucc}.  The generic deduction file
provides substitution, weakening, simultaneous cut, implication and universal
rules, and Leibniz transport.  In particular, the following operations are
already checked:
\[
\begin{array}{ll}
\code{Prf.subst}, & \text{substitute terms for all free variables},\\
\code{Prf.weaken},& \text{enlarge a context},\\
\code{Prf.trans}, & \text{replace hypotheses by their derivations},\\
\code{Prf.eqSubst1},& \text{transport a formula along an equality},\\
\code{PRA.qfInd},& \text{induct on a quantifier-free formula}.
\end{array}
\]
Finally, \code{NativeBiInterpretation.lean} proves
\code{pra_iff_native}: an ordinary iPRA derivation is equivalent to a
derivation of its reading in the native scalar calculus.  Consequently no new
bi-interpretation theorem is needed.  The missing work is entirely on the
ordinary \(\Prf\,\PRAAx\) side.

\subsection{Classification of the reported difficulties}

The difficulty reports in the packets can be separated into four classes.
This classification is essential because only the first class can be solved
by ordinary Lean engineering.

\begin{definition}[Four-way status classification]
A target or subtarget has one of the following statuses.
\begin{enumerate}
\item \emph{Compiler obligation}: all mathematical data and local proofs are
available, but the generic construction of an object-language derivation has
not been implemented.
\item \emph{Quantitative producer obligation}: a finite verifier is known,
but an explicit bound, modulus, representative, or finite cover remains to be
proved.
\item \emph{Representation defect}: the requested output is discontinuous or
noncanonical; a relational or guarded statement is required.
\item \emph{Computability barrier}: the unrestricted producer would imply a
nonconstructive principle or solve an undecidable problem.
\end{enumerate}
\end{definition}

Examples of the first class are bounded-fold reflection and finite-trace
induction.  Examples of the second are general Siegel/Faltings height bounds,
representatives on Mordell--Weil covering curves, and effective constants in
certain prime-distribution results.  Exact root labels at collisions and an
axis for the identity rotation are representation defects.  A total primitive-
recursive universal evaluator and a uniform conversion from mere finiteness to
a search bound are barriers.

The framework below solves the compiler obligations, gives a common interface
for producer obligations, repairs representation defects, and records barriers
rather than concealing them behind axioms.

\subsection{Audit conclusion}

The repository does not need another informal assertion that ``all the
relevant computations are finite.''  It needs a small compilation theorem
whose hypotheses are already the kind of local formulas the packets can
produce and whose conclusion is an actual iPRA derivation.  Once that theorem
is in place, the many domain-specific packets can share a single conversion
workflow:
\[
\begin{aligned}
\text{mathematical proof}
&\longrightarrow \text{finite contracted program}\\
&\longrightarrow \text{local }\Prf\,\PRAAx\text{ obligations}\\
&\longrightarrow \text{global }\Prf\,\PRAAx\text{ derivation}.
\end{aligned}
\]
The remaining sections develop this pipeline and prove its mathematical
correctness.
